% ===================== sec_limitations.tex =====================
% NOTE: each paragraph corresponds to an item in claim_audit_report.md or reconcile_questions.md. Delete the paragraph on the radiology arm only if the arm is re-extracted and rerun.
\section{Limitations}\label{sec:limitations}

\paragraph{Radiology arm.} \rad{Of the 1{,}200 radiology identifiers, 1{,}073 refer to one image file, and 113 have distinct content (Section~\ref{sec:data}). For most radiology items the model saw one radiograph and was scored against labels belonging to other images. Every model-level value includes this arm, which holds about 41\% of the claims. We report the main comparisons with and without radiology, and the specialisation result holds in the other three modalities. The radiology values and the correlations that involve radiology cannot be interpreted, and the statement that radiology has the highest \SMHR{} may reflect the defect.}

\paragraph{Reference labels as ground truth.} The reference is a set of labels from public datasets and not a complete reading of each image. A finding that the labels do not record is treated as unknown by the judge but is usually scored as unsupported, which can inflate fabrication. A finding that the labels record is required to appear in the response, which can inflate omission, and the audit showed that the pipeline calls omission more often than the simulated annotators do. The coverage of the labels differs between the classes and the modalities: in histopathology, half of the patches carry a binary tumour label only, so that the \SMHR{} is a proxy for the hallucination rate in that modality, and in dermoscopy the anatomical site and a specific diagnosis are available for about 61\% and 68\% of the images. The CheXpert labels are extracted automatically from reports.

\paragraph{One automated judge.} All values depend on one 7B model, used with 4-bit quantisation, that extracts and classifies the claims. Only the omission counts are independent of it, and they make up between 25\% and 42\% of the weighted error. Four of the twelve evaluated models belong to the same family as the judge, and a systematic preference of the judge for similar text was not assessed. The tests of label noise used random flips, which average out over many claims and represent the favourable case. The comparison with a second judge of a different family has not been made.

\paragraph{No clinician validation.} No clinician annotated the labels, and the class weights were assigned by the authors. The only independent annotation is a set from six simulated annotators on 48 images, which are clustered, and the agreement of the pipeline with them is moderate for fabrication and weak for most other classes (Section~\ref{sec:validation}). The weights determine what the \SMHR{} means, and their stability has been tested, but their clinical validity has not. The values of the metrics are therefore measures of disagreement with the reference labels, and not measures of clinical harm.

\paragraph{Design of the comparison.} The study has three domain-adapted models and nine general-purpose models that differ in size, architecture and date of release, and it contains no base--specialist pair. The conclusion on specialisation is a comparison with the best general-purpose model, and one of the three domain-adapted models, CheXagent, was developed for a single modality. All models were run with 4-bit quantisation on one GPU. Only open models of up to 12 billion parameters were included, so that the results do not apply to larger or proprietary models or to full-precision inference. Three further models were planned and were not evaluated.
% NOTE: add the reason for the dropped models (reconcile_questions.md, B1).

\paragraph{Single run and prompt.} Each cell of the benchmark has one response, sampled at a temperature of 0.1, and one prompt per modality and query. The variation of the \SMHR{} of a single model across decoding settings (\PH{0.022} to \PH{0.077}) exceeds the differences within the leading group, so that its order is not supported. We did not vary the wording of the prompts, and studies of paired models report that the results depend on it \citep{mclaughlin2026knowledge}. The token limits are those of the final configuration of the notebook, and the result files do not record them, so that their application to every response cannot be proved. One model, LLaVA-Med, produced a large share of responses that begin in mid-sentence, and its values are an upper bound.

\paragraph{Sample size and subsets.} The benchmark has 13 strata of 150 to 400 images. The mitigation and calibration experiments use a subset of 500 images and five models, with 21 to 57 images per stratum and 85 fundus images, so that results by stratum are not reported from them. Several effects in the mitigation experiment changed sign or vanished when the radiology images were removed, and only one of them was supported against both baselines. The knowledge base for retrieval was written by the authors and has not been reviewed by a clinician, retrieval used the text of the query only, and contrastive decoding was run with one setting of its parameters. The planned third strategy, fine-tuning, was not carried out.

\paragraph{Statistical procedures.} The main tables use 2{,}000 stratified bootstrap resamples and the later analyses 5{,}000 resamples without stratification, so that the intervals differ slightly. Tests were adjusted within families defined for the audit, and some of the families are small. Rank correlations between twelve models have wide intervals, and a correlation of the specialisation index across modalities cannot reach significance with three domain-adapted models.

\paragraph{Scope.} The benchmark contains two-dimensional images from four modalities and prompts in English. It does not contain computed tomography, magnetic resonance imaging, ultrasound, video or multiple images for one patient, and it does not evaluate the generation of complete reports. It uses public datasets, whose images may differ from the images of a clinical service, and the datasets may have been seen by some models during training.
